% Zumdahl Chapter 8 test -- 20 MCQ + 2 FRQ. 50 min.
\documentclass{apchem-exam}

\apcunitword{Chapter}
\apcunit{8}
\apcunittitle{Bonding: General Concepts}
\apcdoctitle{Chapter Test}
\apcsubtitle{Zumdahl \S8.1--8.13 \textbullet\ 50 minutes \textbullet\ bond energies provided with Section II}

\begin{document}
\apctitleblock

\examsection{I}{Multiple Choice}{20 questions \textbullet\ 30 minutes \textbullet\ 1 point each}{\nocalculator}

% ---- electronegativity and polarity ----------------------------------------
\begin{mcq}
  Electronegativity generally
  \choices
    \choice{increases down a group and across a period}
    \correct{increases across a period and decreases down a group}
    \choice{decreases across a period and down a group}
    \choice{is the same for all nonmetals}
\end{mcq}
\why{Same Coulombic engine as ionization energy: rising $Z_{\text{eff}}$
across, greater distance down.}

\begin{mcq}
  Which bond is the most polar?
  \choices
    \choice{\ce{C-N}}
    \choice{\ce{C-O}}
    \correct{\ce{H-F}}
    \choice{\ce{O-F}}
\end{mcq}
\why{$\Delta$EN $= 1.9$ --- the largest of the four.}
\distractor{D}{picks two very electronegative atoms, but their DIFFERENCE is
small}

\begin{mcq}
  A molecule with polar bonds will be nonpolar overall if
  \choices
    \choice{the bonds are very strong}
    \correct{its geometry causes the bond dipoles to cancel}
    \choice{it contains no lone pairs}
    \choice{it is a gas at room temperature}
\end{mcq}
\why{Symmetry is what cancels the vector sum.}

\begin{mcq}
  Which molecule is polar?
  \choices
    \choice{\ce{CO2}}
    \choice{\ce{BF3}}
    \choice{\ce{CCl4}}
    \correct{\ce{NH3}}
\end{mcq}
\why{The lone pair on N makes it pyramidal, so the dipoles do not cancel.}

% ---- ions and lattice energy -----------------------------------------------
\begin{mcq}
  Compared with its parent atom, a cation is always
  \choices
    \correct{smaller, because the remaining electrons feel the same nuclear
      charge}
    \choice{larger, because it has lost electron repulsion}
    \choice{the same size}
    \choice{smaller, because it has fewer protons}
\end{mcq}
\why{Often an entire outer shell is removed; the proton count never
changes.}
\distractor{D}{ions never differ from their atoms in proton count}

\begin{mcq}
  For an isoelectronic series of ions, size decreases as
  \choices
    \choice{the number of electrons increases}
    \correct{the nuclear charge increases}
    \choice{the ionic charge becomes more negative}
    \choice{the atomic mass decreases}
\end{mcq}
\why{Same electron count, so only $Z$ varies.}

\begin{mcq}
  Which ion is the smallest?
  \choices
    \choice{\ce{O^2-}}
    \choice{\ce{F-}}
    \choice{\ce{Na+}}
    \correct{\ce{Mg^2+}}
\end{mcq}
\why{All are isoelectronic with Ne; magnesium has the most protons.}

\begin{mcq}
  Lattice energy is proportional to
  \choices
    \correct{$Q_1Q_2/r$}
    \choice{$r/Q_1Q_2$}
    \choice{$Q_1Q_2 \times r$}
    \choice{$Q_1 + Q_2$}
\end{mcq}
\why{A modified Coulomb's law.}

\begin{mcq}
  Which compound has the largest lattice energy magnitude?
  \choices
    \choice{\ce{NaCl}}
    \choice{\ce{KBr}}
    \correct{\ce{MgO}}
    \choice{\ce{CaS}}
\end{mcq}
\why{Charge product 4 AND the smallest such ion pair --- beats \ce{CaS} on
radius.}

\begin{mcq}
  \ce{LiF} and \ce{KI} have the same charge product, but \ce{LiF} has a much
  larger lattice energy because
  \choices
    \correct{its ions are smaller, so they sit closer together}
    \choice{fluorine is more electronegative}
    \choice{lithium has fewer electrons}
    \choice{\ce{KI} is covalent}
\end{mcq}
\why{Smaller $r$ gives stronger Coulombic attraction.}
\distractor{B}{electronegativity governs sharing, not lattice separation}

% ---- bond energies ---------------------------------------------------------
\begin{mcq}
  Breaking a chemical bond is
  \choices
    \choice{always exothermic}
    \correct{always endothermic}
    \choice{exothermic for ionic bonds only}
    \choice{energy-neutral}
\end{mcq}
\why{Energy must be supplied to separate atoms from a lower-energy bonded
state.}

\begin{mcq}
  As bond order increases from single to triple, bond length and bond energy
  respectively
  \choices
    \choice{increase and increase}
    \correct{decrease and increase}
    \choice{increase and decrease}
    \choice{decrease and decrease}
\end{mcq}
\why{More shared pairs pull the nuclei closer and hold them more tightly.}

\begin{mcq}
  $\Delta H$ for a reaction is estimated from bond energies as
  \choices
    \choice{bonds formed minus bonds broken}
    \correct{bonds broken minus bonds formed}
    \choice{the sum of all bond energies}
    \choice{products minus reactants}
\end{mcq}
\why{Note this is the reverse order from the formation-enthalpy equation.}
\distractor{D}{that phrasing belongs to $\Delta H^\circ_f$ calculations}

\begin{mcq}
  Bond-energy calculations give only approximate values because tabulated
  bond energies are
  \choices
    \choice{measured at absolute zero}
    \correct{averages taken over many different compounds}
    \choice{always too large}
    \choice{valid only for ionic substances}
\end{mcq}
\why{A C--H bond in methane differs slightly from one in ethanol.}

% ---- Lewis structures ------------------------------------------------------
\begin{mcq}
  The total number of valence electrons in \ce{ClO3-} is
  \choices
    \choice{24}
    \correct{26}
    \choice{25}
    \choice{32}
\end{mcq}
\why{$7 + 18 + 1 = 26$; the negative charge adds one electron.}
\distractor{A}{forgot the charge}

\begin{mcq}
  Which species has a central atom with an expanded octet?
  \choices
    \choice{\ce{NH3}}
    \choice{\ce{CO2}}
    \correct{\ce{SF6}}
    \choice{\ce{H2O}}
\end{mcq}
\why{Sulfur is in period 3 and holds twelve electrons.}

\begin{mcq}
  \ce{NF5} does not exist because nitrogen
  \choices
    \choice{is too electronegative}
    \correct{is in period 2 and cannot exceed an octet}
    \choice{has too few valence electrons}
    \choice{forms only ionic compounds}
\end{mcq}
\why{Only period 3 and beyond can expand.}

% ---- resonance, formal charge, VSEPR ---------------------------------------
\begin{mcq}
  In \ce{CO3^2-}, the carbon--oxygen bond order is
  \choices
    \choice{1}
    \correct{$4/3$}
    \choice{$3/2$}
    \choice{2}
\end{mcq}
\why{Four bonding pairs spread over three equivalent positions.}
\distractor{C}{that is \ce{NO2-} or \ce{O3}, which have two positions}

\begin{mcq}
  The best experimental evidence that \ce{NO3-} is a resonance hybrid is
  that
  \choices
    \choice{it carries a negative charge}
    \correct{all three N--O bonds have identical, intermediate lengths}
    \choice{nitrogen has a formal charge of $+1$}
    \choice{it is trigonal planar}
\end{mcq}
\why{Alternating structures would produce two different bond lengths.}

\begin{mcq}
  The bond angle in \ce{H2O} is smaller than in \ce{CH4} because water
  \choices
    \choice{has fewer electron domains}
    \correct{has two lone pairs, which repel more strongly than bonding
      pairs}
    \choice{has weaker bonds}
    \choice{is a smaller molecule}
\end{mcq}
\why{Both have four domains; lone pairs compress the angle.}
\distractor{A}{both have four domains --- that is the point}

\examsection{II}{Free Response}{2 questions \textbullet\ 20 minutes}{\calculatorok}

\begin{apcnote}[Bond energies (kJ/mol)]
\ce{H-H} 432 \textbullet\ \ce{Cl-Cl} 239 \textbullet\ \ce{H-Cl} 427
\textbullet\ \ce{N#N} 941 \textbullet\ \ce{N-H} 391 \textbullet\
\ce{O=O} 495 \textbullet\ \ce{O-H} 467 \textbullet\ \ce{C-H} 413
\textbullet\ \ce{C=C} 614 \textbullet\ \ce{C-C} 347
\end{apcnote}

% ---- FRQ 1 -----------------------------------------------------------------
\frq{short}{4}{Zumdahl \S8.8 \textbullet\ bond energies}

Consider \ce{2 H2(g) + O2(g) -> 2 H2O(g)}.

\begin{parts}
  \item Calculate $\Delta H$ using bond energies. Show which bonds are
        broken and which are formed. \workspace[2.6cm]
        \keyonly{\begin{apcanswer}Broken: $2(432) + 495 = 1359$.
        Formed: $4(467) = 1868$. $\Delta H = 1359 - 1868 =
        \SI{-509}{\kilo\joule}$\end{apcanswer}}
  \item Explain why the reaction is exothermic, referring to the relative
        strengths of the bonds involved. \answerlines[2]{The four O--H
        bonds formed release more energy (1868 kJ) than is required to break
        the reactant bonds (1359 kJ), so there is a net release.}
  \item The accepted value is \SI{-483.6}{\kilo\joule}. Explain the
        difference. \answerlines[2]{Tabulated bond energies are averages
        across many compounds, so the method gives an estimate rather than
        an exact value.}
\end{parts}

\begin{rubric}
  \rpoint{1}{(a) correct bonds identified with coefficients applied}
  \rpoint{1}{(a) $-509$ kJ, using broken minus formed --- a positive answer
    earns 0}
  \rpoint{1}{(b) compares total energy released to total required}
  \rpoint{1}{(c) attributes the gap to averaged bond energies}
\end{rubric}

% ---- FRQ 2 -----------------------------------------------------------------
\frq{short}{4}{Zumdahl \S8.4--8.5 \textbullet\ ions and lattice energy}

Consider the compounds \ce{NaF}, \ce{MgO}, and \ce{KCl}.

\begin{parts}
  \item Rank the three by lattice energy magnitude, largest first, and
        justify each comparison. \answerlines[3]{\ce{MgO} $>$ \ce{NaF} $>$
        \ce{KCl}. \ce{MgO} has a charge product of 4 versus 1 for the
        others, which dominates. \ce{NaF} beats \ce{KCl} because \ce{Na+}
        and \ce{F-} are smaller than \ce{K+} and \ce{Cl-}, so $r$ is
        shorter.}
  \item \ce{MgO} melts near \SI{2850}{\celsius} while \ce{NaF} melts near
        \SI{990}{\celsius}. Explain the difference using Coulomb's law.
        \answerlines[2]{Lattice energy scales as $Q_1Q_2/r$. With similar
        ionic radii, \ce{MgO}'s $2{+}/2{-}$ charges give roughly four times
        the attraction, so far more energy is needed to disrupt the
        lattice.}
  \item Rank \ce{Na+}, \ce{F-}, and \ce{Mg^2+} by size and explain the
        principle you used. \answerlines[3]{\ce{F-} $>$ \ce{Na+} $>$
        \ce{Mg^2+}. All three are isoelectronic with neon, so size is
        determined solely by nuclear charge --- more protons pulling the
        same ten electrons gives a smaller ion.}
\end{parts}

\begin{rubric}
  \rpoint{1}{(a) correct order}
  \rpoint{1}{(a) charge product cited BEFORE radius --- radius-first
    reasoning earns 0}
  \rpoint{1}{(b) quantitative Coulombic comparison}
  \rpoint{1}{(c) correct order AND the isoelectronic principle named}
\end{rubric}

\examtotal

\end{document}
